\chapter{Introduction}\label{chap:intro}

Distributed storage services are systems that distribute data across multiple physical nodes to ensure reliable and efficient access even in the event of hardware failures, network issues, or other disruptions. Typically, data is replicated across nodes in diverse geographical locations or data centres, which enhances scalability, performance, and overall reliability. Clients can access these systems through various interfaces, such as REST APIs or file system protocols, allowing them to store and retrieve data seamlessly from around the world. Many of these systems also incorporate mechanisms to maintain consistency during concurrent updates and to detect potential conflicts and some leave this to the application level.

Various types of distributed storage systems have been developed to address different use cases. For example, cloud-based services such as Amazon S3~\cite{amazon_s3} and Google Cloud Storage~\cite{google_cloud_storage} provide scalable storage for large volumes of unstructured data on a subscription basis. In contrast, open source solutions like Ceph~\cite{ceph} and GlusterFS~\cite{glusterfs} can be deployed on-premises or within private clouds. Additionally, consumer-focused platforms such as Dropbox~\cite{dropbox}, Google Drive~\cite{google_drive}, and OneDrive~\cite{onedrive} offer file synchronization and sharing capabilities for real-time collaboration. Furthermore, internal systems like Meta's Tectonic~\cite{tectonic} and Google's Colossus~\cite{colossus}, which function as file systems, are also considered distributed storage systems. While these systems share the common goal of efficient data storage and retrieval, they differ in design objectives and trade-offs. For instance, services like Amazon S3 and Google Cloud Storage are engineered for high availability and durability, whereas internal systems such as Tectonic and Colossus prioritize performance and scalability, often accepting weaker consistency guarantees because such concerns are managed at the application level.

In \Cref{chap:features} (by Bratin) we describe the typical features and use cases of distributed storage systems. In \Cref{chap:design} (by Soukhin) we discuss the basic design issues of these systems, while in \Cref{chap:architecture} (by Swarnabh) we explore the fundamental architectures for their implementation. We then study practical implementations in the subsequent chapters: in \Cref{chap:tectonic} (by Bratin) we examine Tectonic, a distributed storage system developed by Meta, and in \Cref{chap:dropbox} (by Swarnabh) we analyze Dropbox, a widely used consumer-oriented service. In \Cref{chap:ceph} (by Soukhin) we investigate Ceph, an open-source distributed storage system that can be deployed on-premises or in private clouds. Each of these systems has its own unique design objectives and trade-offs, which we will explore in detail.

Finally, in \Cref{chap:conclusion} we summarize our findings and present a comparative analysis of different distributed storage systems, discussing how their design choices impact performance, scalability and guarantees offered to clients.
